\chapter{Stein's Method}
\begin{center}\itshape
Characteristic functions prove limit theorems beautifully, but they need independence and they hide the error.\\
Stein found a way to measure the distance to a Gaussian, or to a Poisson law, with one integration by parts,\\
and it works just as well for dependent pieces: colliding hashes, triangles in random graphs, samples without replacement.
\end{center}

Charles Stein (1920--2016) was a statistician who liked to prove that the obvious thing is wrong. In 1956 he showed that when you estimate three or more
unrelated means from Gaussian data, the obvious estimator, the vector of sample means, can be beaten uniformly: ``Stein's paradox'', which statisticians
still argue about over coffee. Fourteen years later, at the Sixth Berkeley Symposium of 1970, he did the same to the central limit theorem. Everybody
proved it with characteristic functions; Stein proved it with a single identity satisfied by the Gaussian law and by nothing else,
\[
\mathbb E f'(Z)=\mathbb E\,Zf(Z),
\]
and he showed that when a random variable $W$ \emph{almost} satisfies it, $W$ is \emph{almost} Gaussian, with an explicit error. The paper, published in 1972,
treated sums of dependent variables, exactly where Fourier methods struggle.

His student Louis Chen, in a Stanford thesis completed in 1972 and a paper of 1975, ran the same idea for the Poisson law. For a decade few people
noticed; then in 1989 Richard Arratia, Larry Goldstein and Louis Gordon turned the Chen--Stein method into a tool anybody could use, under the title
``Two moments suffice for Poisson approximations''. Today it is how one proves that the number of hash collisions, of near-birthdays, of triangles in a
sparse random graph or of long runs of heads is Poisson, with a bound on the error.

This chapter builds Stein's method twice, for the normal law and for the Poisson law, and uses it on samples without replacement, birthdays and
hashes, and random graphs.

\begin{reminder}[from Volume 3 and Chapters 3 and 4]
The total variation distance between laws on a countable set is $\|\mu-\nu\|=\max_A|\mu(A)-\nu(A)|$ (Volume~3, Chapter~8), and Le Cam's inequality says that a
sum $S$ of independent Bernoulli$(p_i)$ variables with $\sum p_i=\lambda$ satisfies $\|\mathcal L(S)-\text{Poisson}(\lambda)\|\le\sum p_i^2$ (Volume~3, Theorem~9.11).
We write $\varphi(x)=\frac1{\sqrt{2\pi}}e^{-x^2/2}$ and $\Phi$ for the standard Gaussian density and distribution function, and use the lower tail bound
of~\eqref{eq:gausstail}: $1-\Phi(x)\ge\frac x{1+x^2}\varphi(x)$ for $x>0$. By symmetry, $\Phi(x)\ge\frac{|x|}{1+x^2}\varphi(x)$ for $x<0$.
\end{reminder}

\section{The Gaussian integration by parts}

\begin{lemma}[Stein's lemma]\label{lem:stein}
Let $W$ be a random variable with $\mathbb E|W|<\infty$.
\begin{enumerate}[label=(\roman*)]
\item If $W\sim N(0,1)$, then $\mathbb Ef'(W)=\mathbb E\,Wf(W)$ for every $C^1$ function $f$ with $f'$ bounded.
\item Conversely, if $\mathbb Ef'(W)=\mathbb E\,Wf(W)$ for every bounded $C^1$ function $f$ with $f'$ bounded, then $W\sim N(0,1)$.
\end{enumerate}
\end{lemma}
\begin{proof}
(i) Since $\varphi'(x)=-x\varphi(x)$, integration by parts on $[-R,R]$ gives
\[
\int_{-R}^Rf'(x)\varphi(x)\dd x=f(R)\varphi(R)-f(-R)\varphi(-R)+\int_{-R}^Rxf(x)\varphi(x)\dd x.
\]
As $f'$ is bounded, $|f(x)|\le|f(0)|+\|f'\|\,|x|$, so the boundary terms tend to $0$ and both integrals converge absolutely as $R\to\infty$.

(ii) is a consequence of Lemma~\ref{lem:steinsol} below: for $h(x)=\cos tx$ or $h(x)=\sin tx$, that lemma produces a bounded $C^1$ function $f$ with bounded
derivative such that $f'(w)-wf(w)=h(w)-\mathbb Eh(Z)$ for all $w$, where $Z\sim N(0,1)$. Taking expectations at $w=W$, $\mathbb Eh(W)=\mathbb Eh(Z)$. So $W$ and $Z$
have the same characteristic function, hence the same law.
\end{proof}

Part (ii) is the seed of the method. Fix a test function $h$ and solve the \emph{Stein equation}
\begin{equation}\label{eq:steineq}
f'(w)-wf(w)=h(w)-\mathbb Eh(Z).
\end{equation}
Put $W$ in it and take expectations:
\begin{equation}\label{eq:steinid}
\mathbb Eh(W)-\mathbb Eh(Z)=\mathbb E\bigl[f'(W)-Wf(W)\bigr].
\end{equation}
The left side is what we want to bound, a question about two laws. The right side involves only $W$, and if $f$ is smooth with controlled derivatives we
can attack it with Taylor's formula, independence, symmetry, whatever structure $W$ has. Everything rests on bounds for the solution~$f$.

\begin{think}
Why not simply check that $W$ has the Gaussian moments $\mathbb EW^k$? The identity of Lemma~\ref{lem:stein} with $f(w)=w^{k-1}$ gives
$\mathbb EZ^k=(k-1)\mathbb EZ^{k-2}$, and the moments do characterize the Gaussian. What does the method of moments \emph{not} give you, that a bound on
$\mathbb E[f'(W)-Wf(W)]$ for a whole class of $f$ will?
\end{think}

\section{Solving the Stein equation}

\begin{lemma}[the solution and its derivatives]\label{lem:steinsol}
Let $h:\R\to\R$ be $C^1$ with $\|h'\|=\sup|h'|<\infty$, and $\tilde h=h-\mathbb Eh(Z)$. The function
\[
f(w)=\frac1{\varphi(w)}\int_{-\infty}^w\tilde h(x)\varphi(x)\dd x
\]
is $C^2$, solves~\eqref{eq:steineq}, and satisfies
\[
\|f\|\le\|h'\|,\qquad\|f'\|\le2\|h'\|,\qquad\|f''\|\le2\|h'\|.
\]
\end{lemma}
\begin{proof}
Since $\frac{\dd}{\dd w}\frac1{\varphi(w)}=\frac w{\varphi(w)}$, differentiating gives $f'(w)=wf(w)+\tilde h(w)$, which is~\eqref{eq:steineq}; $f'$ is continuous,
and so is $f''=f+wf'+h'$. The work is in the bounds, and the trick is to write everything in terms of $h'$.

\emph{Step 1: $\tilde h$ in terms of $h'$.} Writing $h(x)-h(y)=\int_y^xh'(z)\dd z$ and integrating against $\varphi(y)\dd y$, Fubini gives
\[
\tilde h(x)=\int_{\R}\bigl(h(x)-h(y)\bigr)\varphi(y)\dd y=\int_{-\infty}^xh'(z)\Phi(z)\dd z-\int_x^\infty h'(z)\bigl(1-\Phi(z)\bigr)\dd z,
\]
since $z<x$ lies between $y$ and $x$ exactly when $y<z$, which has weight $\Phi(z)$, and $z>x$ lies between $x$ and $y$ when $y>z$, of weight $1-\Phi(z)$.
Call the two integrals $A(x)$ and $B(x)$, so $\tilde h=A-B$.

\emph{Step 2: $f$ in terms of $h'$.} Integrate $\tilde h\varphi$ up to $w$ and use Fubini again: $\int_{-\infty}^wA\varphi=\int_{-\infty}^wh'(z)\Phi(z)\bigl(\Phi(w)-\Phi(z)\bigr)\dd z$
and $\int_{-\infty}^wB\varphi=\int_{-\infty}^wh'(z)(1-\Phi(z))\Phi(z)\dd z+\Phi(w)B(w)$. Subtracting,
\begin{equation}\label{eq:steinrep}
f(w)=-\frac{\bigl(1-\Phi(w)\bigr)A(w)+\Phi(w)B(w)}{\varphi(w)}.
\end{equation}
Introduce the two positive functions
\[
\begin{gathered}
\psi(w)=\int_w^\infty\bigl(1-\Phi(z)\bigr)\dd z=\varphi(w)-w\bigl(1-\Phi(w)\bigr),\\\psi(-w)=\int_{-\infty}^w\Phi(z)\dd z=\varphi(w)+w\Phi(w)
\end{gathered}
\]
(differentiate the right sides, using $\varphi'=-w\varphi$, and check the limits). Then $|A(w)|\le\|h'\|\psi(-w)$ and $|B(w)|\le\|h'\|\psi(w)$.

\emph{Step 3: $\|f\|$.} By~\eqref{eq:steinrep}, $|f(w)|\le\|h'\|\bigl((1-\Phi)\psi(-w)+\Phi\psi(w)\bigr)/\varphi$, and
\[
(1-\Phi)(\varphi+w\Phi)+\Phi\bigl(\varphi-w(1-\Phi)\bigr)=\varphi.
\]
So $|f|\le\|h'\|$.

\emph{Step 4: $\|f'\|$.} From $f'=wf+A-B$ and~\eqref{eq:steinrep},
\[
f'(w)=A(w)\Bigl(1-\frac{w(1-\Phi(w))}{\varphi(w)}\Bigr)-B(w)\Bigl(1+\frac{w\Phi(w)}{\varphi(w)}\Bigr)=\frac{A(w)\psi(w)-B(w)\psi(-w)}{\varphi(w)}.
\]
Hence $|f'(w)|\le\|h'\|\cdot2\psi(w)\psi(-w)/\varphi(w)$, an even function of $w$. For $w\ge0$ the tail bound gives
$\psi(w)\le\varphi(w)-\frac{w^2}{1+w^2}\varphi(w)=\frac{\varphi(w)}{1+w^2}$, and $\psi(-w)\le\varphi(0)+w$, so
\[
\frac{2\psi(w)\psi(-w)}{\varphi(w)}\le\frac{2\varphi(0)}{1+w^2}+\frac{2w}{1+w^2}\le2\varphi(0)+1<2.
\]

\emph{Step 5: $\|f''\|$.} From $f''=f+wf'+h'=(1+w^2)f+w\tilde h+h'$ and~\eqref{eq:steinrep},
\[
\begin{gathered}
(1+w^2)f+w\tilde h=A(w)\,\alpha(w)+B(w)\,\beta(w),\\
\alpha=w-\frac{(1+w^2)(1-\Phi)}{\varphi},\qquad\beta=-w-\frac{(1+w^2)\Phi}\varphi.
\end{gathered}
\]
The tail bound says exactly that $\alpha\le0$ for $w>0$, and $\alpha\le0$ is clear for $w\le0$; likewise $\beta\le0$ everywhere. So
\[
\bigl|(1+w^2)f+w\tilde h\bigr|\le\|h'\|\bigl(-\alpha\,\psi(-w)-\beta\,\psi(w)\bigr),
\]
and expanding with the formulas for $\psi(\pm w)$, the $w\varphi$ terms cancel, the $\Phi(1-\Phi)$ terms cancel, and what is left is
\[
-w(\varphi+w\Phi)+w\bigl(\varphi-w(1-\Phi)\bigr)+\frac{1+w^2}{\varphi}\bigl((1-\Phi)\varphi+\Phi\varphi\bigr)=-w^2+(1+w^2)=1.
\]
Therefore $|f''|\le\|h'\|+|h'|\le2\|h'\|$.
\end{proof}

The constant $2$ for $f'$ is not sharp: the function $2\psi(w)\psi(-w)/\varphi(w)$ has its maximum $2\varphi(0)=\sqrt{2/\pi}\approx0.80$ at $w=0$, a fact proved
in Chen, Goldstein and Shao's book (Lemma~2.4) and checked numerically. The constant $1$ for $f$ is sharp: $h(x)=x$ gives $f\equiv-1$.

\begin{figure}[ht]
\centering\fig{ch06-steinsol}
\caption{The solution $f$ of the Stein equation for $h(x)=|x|$, and its derivative. The dashed envelope $\pm2\psi(w)\psi(-w)/\varphi(w)$ bounds $|f'|$ for every
$1$-Lipschitz $h$; here $f'$ touches it at $w=0$, where it equals $\sqrt{2/\pi}$. (This $h$ has a corner at $0$; Steps~1--4 of the proof only use $|h'|\le1$ almost everywhere, so they apply.)}
\label{fig:steinsol}
\end{figure}

\section{Wasserstein distance and the central limit theorem}

Lemma~\ref{lem:steinsol} controls $f$ through $\|h'\|$, so the natural distance is the one tested by Lipschitz functions. For random variables $X,Y$ with
finite means, the \emph{Wasserstein distance} is
\[
d_W(X,Y)=\sup\bigl\{|\mathbb Eh(X)-\mathbb Eh(Y)|:\ |h(x)-h(y)|\le|x-y|\text{ for all }x,y\bigr\}.
\]
It depends only on the laws, and $d_W(X_n,Y)\to0$ implies $X_n\Rightarrow Y$ (Exercise~\ref{ex:dwlaw}). Two remarks make it usable. First, the supremum may be taken
over $C^1$ functions with $\|h'\|\le1$: if $h$ is $1$-Lipschitz and $N\sim N(0,1)$, then $h_\varepsilon(x)=\mathbb Eh(x+\varepsilon N)$ is $C^1$ (differentiate the Gaussian
density), $1$-Lipschitz, and $|h_\varepsilon-h|\le\varepsilon\mathbb E|N|\le\varepsilon$. Second, it controls the Kolmogorov distance
$d_K(X,Y)=\sup_x|P(X\le x)-P(Y\le x)|$ to a Gaussian.

\begin{proposition}\label{prop:dkdw}
For any $W$ with finite mean,
\[
d_K(W,Z)\le\sqrt{2\varphi(0)\,d_W(W,Z)}=(2/\pi)^{1/4}\sqrt{d_W(W,Z)}.
\]
\end{proposition}
\begin{proof}
Fix $x$ and $\varepsilon>0$, and let $h$ be $1$ on $(-\infty,x]$, $0$ on $[x+\varepsilon,\infty)$ and linear in between. Then $\varepsilon h$ is $1$-Lipschitz and
\[
\begin{gathered}
P(W\le x)-\Phi(x)\le\mathbb Eh(W)-\mathbb Eh(Z)+\mathbb Eh(Z)-\Phi(x)\\\le\frac{d_W}\varepsilon+\int_x^{x+\varepsilon}\Bigl(1-\frac{t-x}\varepsilon\Bigr)\varphi(t)\dd t
\le\frac{d_W}\varepsilon+\frac{\varepsilon\varphi(0)}2.
\end{gathered}
\]
The ramp from $x-\varepsilon$ to $x$ gives the same bound for $\Phi(x)-P(W\le x)$. Take $\varepsilon=\sqrt{2d_W/\varphi(0)}$.
\end{proof}

\begin{theorem}[CLT with an error in Wasserstein distance]\label{thm:steinclt}
Let $X_1,\dots,X_n$ be independent, with $\mathbb EX_i=0$, $\sum_i\mathbb EX_i^2=1$ and $\mathbb E|X_i|^3<\infty$, and let $W=X_1+\dots+X_n$. Then
\[
d_W(W,Z)\le3\sum_{i=1}^n\mathbb E|X_i|^3.
\]
In particular, for i.i.d.\ $\xi_i$ with mean $0$, variance $\sigma^2$ and $\rho=\mathbb E|\xi|^3/\sigma^3$,
$d_W\bigl(\frac{\xi_1+\dots+\xi_n}{\sigma\sqrt n},Z\bigr)\le\frac{3\rho}{\sqrt n}$.
\end{theorem}
\begin{proof}
Let $h$ be $C^1$ with $\|h'\|\le1$ and $f$ the solution of Lemma~\ref{lem:steinsol}; by~\eqref{eq:steinid} we must bound $\mathbb E[f'(W)-Wf(W)]$. Let
$\sigma_i^2=\mathbb EX_i^2$ and $W^{(i)}=W-X_i$, which is independent of $X_i$. Since $\mathbb EX_i=0$,
\[
\mathbb E\,X_if(W)=\mathbb E\,X_i\bigl(f(W)-f(W^{(i)})\bigr)=\mathbb E\Bigl[X_i^2\int_0^1f'\bigl(W^{(i)}+tX_i\bigr)\dd t\Bigr].
\]
Sum over $i$, and write $\mathbb Ef'(W)=\sum_i\sigma_i^2\mathbb Ef'(W)$ and $\mathbb E\,X_i^2f'(W^{(i)})=\sigma_i^2\mathbb Ef'(W^{(i)})$ by independence:
\[
\begin{aligned}
\mathbb E\bigl[f'(W)-Wf(W)\bigr]&=\sum_i\sigma_i^2\,\mathbb E\bigl[f'(W)-f'(W^{(i)})\bigr]\\
&\quad-\sum_i\mathbb E\Bigl[X_i^2\int_0^1\bigl(f'(W^{(i)}+tX_i)-f'(W^{(i)})\bigr)\dd t\Bigr].
\end{aligned}
\]
The first sum is at most $\|f''\|\sum_i\sigma_i^2\mathbb E|X_i|$ in absolute value, and the second at most $\|f''\|\sum_i\mathbb E|X_i|^3\int_0^1t\dd t$. By
Jensen's inequality $\sigma_i^2\le(\mathbb E|X_i|^3)^{2/3}$ and $\mathbb E|X_i|\le(\mathbb E|X_i|^3)^{1/3}$, so $\sigma_i^2\mathbb E|X_i|\le\mathbb E|X_i|^3$. With $\|f''\|\le2$
the total is at most $2\sum\mathbb E|X_i|^3+\sum\mathbb E|X_i|^3$. The i.i.d.\ case is $X_i=\xi_i/(\sigma\sqrt n)$.
\end{proof}

The order $n^{-1/2}$ is right: for fair $\pm1$ steps the lattice forces $d_W\ge\frac1{2\sqrt n}-O(\frac1n)$ (Exercise~\ref{ex:latticedw}). Through
Proposition~\ref{prop:dkdw} the theorem only gives $d_K=O(n^{-1/4})$, worse than Berry--Esseen (Theorem~\ref{thm:berryesseen}); Stein's method can recover
$n^{-1/2}$ for $d_K$ too, with a ``concentration inequality'' that is more delicate (Chen, Goldstein and Shao, Chapter~3). The point of the proof above is
different: it used independence in exactly one place, to write $\mathbb EX_if(W^{(i)})=0$ and $\mathbb EX_i^2f'(W^{(i)})=\sigma_i^2\mathbb Ef'(W^{(i)})$. Any structure
that gives approximate versions of these two facts gives a CLT with an error bound.

\section{Exchangeable pairs}

Stein's own device for dependence: build a second copy $W'$ of $W$ by a small random change, so that the pair $(W,W')$ is \emph{exchangeable},
$(W,W')\overset{\mathrm d}=(W',W)$.

\begin{theorem}[exchangeable pairs]\label{thm:exch}
Let $(W,W')$ be exchangeable with $\mathbb EW^2=1$, let $\mathcal G$ be a $\sigma$-field for which $W$ is measurable, and suppose that for some $\lambda>0$
\[
\mathbb E[W'-W\mid\mathcal G]=-\lambda W.
\]
Then, with $D=W'-W$,
\[
d_W(W,Z)\le2\,\mathbb E\Bigl|1-\frac{\mathbb E[D^2\mid\mathcal G]}{2\lambda}\Bigr|+\frac{\mathbb E|D|^3}{2\lambda}.
\]
\end{theorem}
\begin{proof}
Let $f$ be bounded, $C^2$, with $f',f''$ bounded. By exchangeability, the antisymmetric function $(w,w')\mapsto(w'-w)(f(w')+f(w))$ has mean $0$ at $(W,W')$.
Writing $f(W')+f(W)=(f(W')-f(W))+2f(W)$ and using $\mathbb E[Df(W)]=\mathbb E\bigl[f(W)\mathbb E[D\mid\mathcal G]\bigr]=-\lambda\mathbb E\,Wf(W)$,
\[
\mathbb E\,Wf(W)=\frac1{2\lambda}\mathbb E\bigl[D\bigl(f(W')-f(W)\bigr)\bigr]=\frac1{2\lambda}\mathbb E\bigl[D^2f'(W)\bigr]+\frac1{2\lambda}\mathbb E[DR],
\]
where $R=f(W')-f(W)-Df'(W)$ satisfies $|R|\le\frac12\|f''\|D^2$. Conditioning the first term on $\mathcal G$,
\[
\mathbb E\bigl[f'(W)-Wf(W)\bigr]=\mathbb E\Bigl[f'(W)\Bigl(1-\frac{\mathbb E[D^2\mid\mathcal G]}{2\lambda}\Bigr)\Bigr]-\frac1{2\lambda}\mathbb E[DR],
\]
which is at most $\|f'\|\,\mathbb E|1-\mathbb E[D^2\mid\mathcal G]/(2\lambda)|+\|f''\|\,\mathbb E|D|^3/(4\lambda)$ in absolute value. Apply this to the solution of Lemma~\ref{lem:steinsol}
for $C^1$ functions $h$ with $\|h'\|\le1$, so that $\|f'\|,\|f''\|\le2$. (All expectations are finite if $\mathbb E|D|^3<\infty$; otherwise there is nothing to prove.)
\end{proof}

Two sanity checks. Taking expectations in the hypothesis gives $\mathbb EW=0$; and $\mathbb E\,WW'=(1-\lambda)\mathbb EW^2$, so $\mathbb ED^2=2\lambda$: the quantity
$\mathbb E[D^2\mid\mathcal G]/(2\lambda)$ always has mean $1$, and the first term of the bound measures how much it fluctuates.

\begin{example}[sampling without replacement]\label{ex:hypergeom}
An urn holds $N$ balls, $K$ of them white; we draw $m$ without replacement and count the white ones, $S$. This is the situation of every opinion poll, and of
every audit that checks $m$ invoices out of $N$. With $p=K/N$, $\mathbb ES=mp$ and $\sigma^2=\operatorname{Var}S=mp(1-p)\frac{N-m}{N-1}$. Then
\[
d_W\Bigl(\frac{S-mp}\sigma,Z\Bigr)\le\frac2\sigma+\frac4N.
\]
\end{example}
\begin{proof}
Make the exchangeable pair by swapping a uniformly chosen ball of the sample with a uniformly chosen ball outside it. This move from one $m$-subset to
another has the same probability $\frac1{m(N-m)}$ as the reverse move, and the sample is uniform, so (sample, new sample) is exchangeable, and so is $(S,S')$.
The count changes by $+1$ if we discard a black ball and bring in a white one, by $-1$ in the opposite case:
\[
P(S'-S=1\mid S)=\frac{(m-S)(K-S)}{m(N-m)},\qquad P(S'-S=-1\mid S)=\frac{S(N-m-K+S)}{m(N-m)}.
\]
Subtracting, the $S^2$ and $mS$ and $KS$ terms cancel: $\mathbb E[S'-S\mid S]=\frac{mK-NS}{m(N-m)}=-\frac N{m(N-m)}(S-mp)$. So
$W=(S-mp)/\sigma$ satisfies the hypothesis with $\mathcal G=\sigma(S)$ and $\lambda=\frac N{m(N-m)}$. Adding instead of subtracting, and writing $S=mp+x$,
\[
\begin{gathered}
m(N-m)\,\mathbb E[(S'-S)^2\mid S]=mK+S(N-2m-2K)+2S^2\\=c+N(1-2a)(1-2p)\,x+2(x^2-\sigma^2),
\end{gathered}
\]
with $a=\frac mN$ and a constant $c$; the constant is whatever makes the mean right, since $\mathbb E[D^2]=2\lambda$. As $D^2=(S'-S)^2/\sigma^2$ and $2\lambda m(N-m)\sigma^2=2N\sigma^2$,
\[
\mathbb E\Bigl|1-\frac{\mathbb E[D^2\mid S]}{2\lambda}\Bigr|=\frac{\mathbb E\bigl|N(1-2a)(1-2p)x+2(x^2-\sigma^2)\bigr|}{2N\sigma^2}\le\frac{N\sigma+4\sigma^2}{2N\sigma^2}=\frac1{2\sigma}+\frac2N,
\]
using $\mathbb E|x|\le\sigma$ and $\mathbb E|x^2-\sigma^2|\le2\sigma^2$. Finally $|D|$ is $0$ or $\frac1\sigma$, so $\mathbb E|D|^3=\frac1\sigma\mathbb ED^2=\frac{2\lambda}\sigma$.
Theorem~\ref{thm:exch} gives $2\bigl(\frac1{2\sigma}+\frac2N\bigr)+\frac1\sigma$.
\end{proof}

The bound depends on $m$ and $N$ only through $\sigma$, and does not care whether $m$ is a small or large fraction of $N$, the case where binomial
approximations fail. For a poll of $m=1000$ out of $N=10^4$ with $p=\frac12$, $\sigma\approx15$ and the bound is $0.13$, while the true distance, computed numerically, is $0.017$. Stein bounds are honest, not tight.

\section{The Chen--Stein method}

The Poisson law has its own integration by parts. If $W\sim$ Poisson$(\lambda)$ and $g$ is bounded on $\N$, then
\[
\mathbb E\,Wg(W)=\sum_{k\ge1}kg(k)e^{-\lambda}\frac{\lambda^k}{k!}=\lambda\sum_{k\ge1}g(k)e^{-\lambda}\frac{\lambda^{k-1}}{(k-1)!}=\lambda\,\mathbb Eg(W+1),
\]
and conversely this identity characterizes the Poisson law (Exercise~\ref{ex:poissonchar}). Write $\pi_k=e^{-\lambda}\frac{\lambda^k}{k!}$, $\mathrm{Po}(A)=\sum_{k\in A}\pi_k$ and
$U_k=\{0,1,\dots,k\}$. The Stein equation for a set $A\subseteq\N$ is
\begin{equation}\label{eq:chensteineq}
\lambda g(k+1)-kg(k)=\mathbf 1_A(k)-\mathrm{Po}(A),\qquad k\ge0.
\end{equation}
Multiply by $\pi_k$ and use $k\pi_k=\lambda\pi_{k-1}$: the left side becomes $\lambda\bigl(\pi_kg(k+1)-\pi_{k-1}g(k)\bigr)$, which telescopes. With $g(0)=0$ the
unique solution is, for $k\ge0$,
\begin{equation}\label{eq:chensteinsol}
g_A(k+1)=\frac{\mathrm{Po}(A\cap U_k)-\mathrm{Po}(A)\,\mathrm{Po}(U_k)}{\lambda\pi_k}.
\end{equation}

\begin{lemma}[the Poisson Stein solution]\label{lem:chensteinsol}
For every $A\subseteq\N$, $\displaystyle\sup_{k\ge0}|g_A(k+1)-g_A(k)|\le\frac{1-e^{-\lambda}}\lambda\le\min\Bigl(1,\frac1\lambda\Bigr)$.
\end{lemma}
\begin{proof}
Write $\Delta g(k)=g(k+1)-g(k)$. The solution is linear in $\mathbf 1_A$, so $g_A=\sum_{j\in A}g_j$ with $g_j=g_{\{j\}}$, and $\sum_{j\ge0}g_j=g_\N=0$ because
the right side of~\eqref{eq:chensteineq} vanishes for $A=\N$.

\emph{Signs.} From~\eqref{eq:chensteinsol}, for $k\ge j$,
\[
g_j(k+1)=\frac{\pi_j\bigl(1-\mathrm{Po}(U_k)\bigr)}{\lambda\pi_k}=\pi_j\sum_{i>k}\frac{k!\,\lambda^{i-k-1}}{i!}=\pi_j\sum_{r\ge0}\frac{k!\,\lambda^r}{(k+1+r)!},
\]
positive and decreasing in $k$; and for $k<j$,
\[
g_j(k+1)=-\frac{\pi_j\mathrm{Po}(U_k)}{\lambda\pi_k}=-\pi_j\sum_{r=0}^k\frac{k(k-1)\cdots(k-r+1)}{\lambda^{r+1}},
\]
negative and decreasing in $k$, as is $g_j(0)=0\ge g_j(1)$. So $g_j$ decreases on $\{0,\dots,j\}$ and on $\{j+1,j+2,\dots\}$: $\Delta g_j(k)\le0$ for every $k\ne j$.

\emph{Reduction.} Fix $k$. If $k\in A$, then $\Delta g_A(k)=\sum_{j\in A}\Delta g_j(k)\le\Delta g_k(k)$, dropping the non-positive terms $j\ne k$; if $k\notin A$,
$\Delta g_A(k)\le0$. Applying this to $A^c$ and using $g_A=-g_{A^c}$ gives the lower bound $\Delta g_A(k)\ge-\Delta g_k(k)$. So $|\Delta g_A(k)|\le\Delta g_k(k)$.

\emph{The value.} By the two formulas, $g_k(k+1)=\frac1\lambda\sum_{i>k}\pi_i$ and $g_k(k)=-\frac1k\sum_{i<k}\pi_i$ (and $g_0(0)=0$). Since $\frac1k\le\frac1{i+1}$ for $i<k$,
\[
\Delta g_k(k)=e^{-\lambda}\Bigl(\sum_{i>k}\frac{\lambda^{i-1}}{i!}+\sum_{i<k}\frac{\lambda^i}{i!\,k}\Bigr)\le e^{-\lambda}\Bigl(\sum_{i>k}\frac{\lambda^{i-1}}{i!}+\sum_{i=1}^{k}\frac{\lambda^{i-1}}{i!}\Bigr)
=e^{-\lambda}\frac{e^\lambda-1}\lambda.
\]
Finally $1-e^{-\lambda}\le\min(1,\lambda)$.
\end{proof}

Now let $W=\sum_{\alpha\in\Gamma}I_\alpha$ be a finite sum of indicators, $p_\alpha=\mathbb EI_\alpha$, $\lambda=\sum p_\alpha=\mathbb EW$. Suppose each $\alpha$ has a
\emph{neighbourhood} $B_\alpha\ni\alpha$ such that $I_\alpha$ is independent of the family $(I_\beta)_{\beta\notin B_\alpha}$. Put
\[
b_1=\sum_{\alpha\in\Gamma}\sum_{\beta\in B_\alpha}p_\alpha p_\beta,\qquad b_2=\sum_{\alpha\in\Gamma}\sum_{\beta\in B_\alpha\setminus\{\alpha\}}\mathbb E[I_\alpha I_\beta].
\]
$b_1$ is small when the neighbourhoods are small; $b_2$ is small when neighbouring events rarely happen together, so that occurrences do not come in clumps.

\begin{theorem}[Chen--Stein, local dependence]\label{thm:chenstein}
Under these assumptions, $\displaystyle\|\mathcal L(W)-\text{Poisson}(\lambda)\|\le(b_1+b_2)\frac{1-e^{-\lambda}}\lambda$.
\end{theorem}
\begin{proof}
Fix $A$ and let $g=g_A$; by~\eqref{eq:chensteineq} at $k=W$, $P(W\in A)-\mathrm{Po}(A)=\mathbb E[\lambda g(W+1)-Wg(W)]=\sum_\alpha\bigl(p_\alpha\mathbb Eg(W+1)-\mathbb E\,I_\alpha g(W)\bigr)$.
Fix $\alpha$ and split $W=I_\alpha+Z_\alpha+V_\alpha$ with $Z_\alpha=\sum_{\beta\in B_\alpha\setminus\{\alpha\}}I_\beta$ and $V_\alpha=\sum_{\beta\notin B_\alpha}I_\beta$, independent of $I_\alpha$.
On $\{I_\alpha=1\}$, $W=1+Z_\alpha+V_\alpha$, so
\[
\mathbb E\,I_\alpha g(W)=\mathbb E\,I_\alpha g(V_\alpha+1)+\mathbb E\bigl[I_\alpha\bigl(g(1+Z_\alpha+V_\alpha)-g(1+V_\alpha)\bigr)\bigr],
\]
and $\mathbb E\,I_\alpha g(V_\alpha+1)=p_\alpha\mathbb Eg(V_\alpha+1)$ by independence. Hence
\[
\begin{aligned}
&p_\alpha\mathbb Eg(W+1)-\mathbb E\,I_\alpha g(W)\\
&\quad=p_\alpha\mathbb E\bigl[g(W+1)-g(V_\alpha+1)\bigr]\\
&\qquad-\mathbb E\bigl[I_\alpha\bigl(g(1+Z_\alpha+V_\alpha)-g(1+V_\alpha)\bigr)\bigr].
\end{aligned}
\]
A difference $g(a+b)-g(a)$ with $b\ge0$ is a sum of $b$ increments, so it is at most $b\sup|\Delta g|$. Here $W-V_\alpha=\sum_{\beta\in B_\alpha}I_\beta$ and the second
difference is at most $Z_\alpha\sup|\Delta g|$, so the right side is at most $\sup|\Delta g|\bigl(\sum_{\beta\in B_\alpha}p_\alpha p_\beta+\mathbb E\,I_\alpha Z_\alpha\bigr)$.
Sum over $\alpha$, use Lemma~\ref{lem:chensteinsol}, and take the supremum over $A$.
\end{proof}

With independent indicators take $B_\alpha=\{\alpha\}$: $b_2=0$, $b_1=\sum p_\alpha^2$, and the bound is $\frac{1-e^{-\lambda}}\lambda\sum p_\alpha^2$, which improves Le Cam's
inequality by the factor $\min(1,\frac1\lambda)$, as promised in Volume~3. For binomial$(n,\frac\lambda n)$ it reads $\frac{\lambda(1-e^{-\lambda})}n\le\frac\lambda n$ instead of $\frac{\lambda^2}n$.

\section{Birthdays, hashes and triangles}

\begin{example}[birthdays and hash collisions]\label{ex:hashcs}
Hash $n$ keys independently and uniformly into $d$ buckets, and let $W$ be the number of colliding pairs. Then $\lambda=\binom n2/d$ and
\[
\Bigl|P(\text{no collision})-e^{-\lambda}\Bigr|\le\|\mathcal L(W)-\text{Poisson}(\lambda)\|\le\frac{(4n-7)(1-e^{-\lambda})}d.
\]
\end{example}
\begin{proof}
Index by pairs $\alpha=\{i,j\}$, $I_\alpha=\mathbf 1\{h(i)=h(j)\}$, $p_\alpha=\frac1d$. $I_{\{i,j\}}$ is a function of $h(i),h(j)$, so it is independent of the
indicators of pairs disjoint from $\{i,j\}$; take $B_\alpha$ to be the $2(n-2)+1=2n-3$ pairs meeting $\alpha$. For $\beta\ne\alpha$ in $B_\alpha$, say $\alpha=\{i,j\}$
and $\beta=\{i,k\}$, $P(h(i)=h(j)=h(k))=\frac1{d^2}$: the collision indicators are pairwise independent even when they overlap. So
$b_1=\binom n2(2n-3)\frac1{d^2}$, $b_2=\binom n2\,2(n-2)\frac1{d^2}$, and $(b_1+b_2)/\lambda=\frac{4n-7}d$.
\end{proof}

For $23$ people and $365$ days, $\lambda=0.693$ and $e^{-\lambda}=0.500$, while the exact probability of no shared birthday is $0.493$ (Volume~3, Proposition~1.8);
the bound $0.12$ is true but loose (Figure~\ref{fig:birthday}). The bound shines where the exact product is unpleasant and the stakes are real. With $32$-bit identifiers
($d=2^{32}$) and $n=10^5$ objects, $\lambda=1.164$: the probability that all identifiers are distinct is $e^{-1.164}=0.312$ up to an error at most
$6.5\times10^{-5}$. With $77\,163$ objects the probability of a collision is $\frac12$, within $4\times10^{-5}$. A $32$-bit identifier is not unique for a hundred
thousand objects, and the Chen--Stein bound says so with a certificate.

\begin{figure}[ht]
\centering\fig{ch06-birthday}
\caption{The birthday problem with $d=365$. Exact probability of no shared birthday among $n$ people (dots), the Poisson approximation $e^{-\binom n2/365}$ (line),
and the Chen--Stein band of half-width $(4n-7)(1-e^{-\lambda})/365$ (shaded). The true error never exceeds $0.011$; the band is wider, but it is a proof.}
\label{fig:birthday}
\end{figure}

\begin{example}[triangles in a sparse random graph]\label{ex:triangles}
In the random graph $G(n,p)$ of Volume~3 (each of the $\binom n2$ edges present independently with probability $p$) with $p=c/n$, let $T$ be the number of
triangles. Then $\|\mathcal L(T)-\text{Poisson}(\lambda_n)\|\le\frac{c^5}{2n}+\frac{c^6}{2n^2}$ with $\lambda_n=\binom n3p^3\to\frac{c^3}6$. In particular
$P(G(n,c/n)\text{ has no triangle})\to e^{-c^3/6}$.
\end{example}
\begin{proof}
Index by triples $\alpha$ of vertices, and let $I_\alpha$ indicate that the three edges of $\alpha$ are present, so $p_\alpha=p^3$. Triples that share at most one vertex share
no edge, so $I_\alpha$ is independent of the indicators of those triples; take $B_\alpha$ to be $\alpha$ and the $3(n-3)$ triples sharing an edge with $\alpha$.
Two triangles sharing an edge use $5$ edges, so
\[
\begin{gathered}b_1=\binom n3\bigl(3n-8\bigr)p^6\le\frac{n^3}6\cdot3n\cdot\frac{c^6}{n^6},\qq b_2=\binom n3\,3(n-3)\,p^5\le\frac{n^3}6\cdot3n\cdot\frac{c^5}{n^5}.\end{gathered}
\]
Since $\frac{1-e^{-\lambda}}\lambda\le1$, Theorem~\ref{thm:chenstein} gives the bound, and $\lambda_n=\frac{n(n-1)(n-2)}6\frac{c^3}{n^3}\to\frac{c^3}6$.
\end{proof}

In the dense regime $p$ fixed, $\lambda\to\infty$ and the triangle count is approximately normal instead; the bound of Theorem~\ref{thm:chenstein} then grows, as it
should, since a Poisson law with large mean is itself close to normal and is no longer the right description of the fluctuations.

\begin{think}
The number of fixed points of a uniform random permutation tends to Poisson$(1)$ (Volume~3, Theorem~2.13). Write it as $\sum_i\mathbf 1\{\sigma(i)=i\}$. Why does
Theorem~\ref{thm:chenstein} not apply directly? The indicators are only weakly dependent, but none is independent of the others. Arratia, Goldstein and
Gordon's paper, and Barbour, Holst and Janson's book, handle such cases with \emph{couplings} instead of neighbourhoods: build, next to $W$, a variable with the
law of $W-1$ given $I_\alpha=1$, and measure how far it is from $W-I_\alpha$.
\end{think}

\section{Exercises}

\begin{exercise}[$\star$]\label{ex:gaussmoments}
Use Lemma~\ref{lem:stein} to compute $\mathbb EZ^{2k}$ for $Z\sim N(0,1)$, and show that $X\sim N(\mu,s^2)$ if and only if $\mathbb E(X-\mu)f(X)=s^2\mathbb Ef'(X)$ for all
bounded $C^1$ functions $f$ with $f'$ bounded.
\end{exercise}

\begin{exercise}[$\star$]\label{ex:dwlaw}
Show that $d_W(X_n,Y)\to0$ implies $X_n\Rightarrow Y$.
\end{exercise}

\begin{exercise}[$\star$]\label{ex:poissonchar}
Let $W$ take values in $\N$ with $\mathbb EW<\infty$. Show that $W\sim$ Poisson$(\lambda)$ if and only if $\mathbb E\,Wg(W)=\lambda\mathbb Eg(W+1)$ for every bounded $g$. Deduce the
factorial moments $\mathbb E\,W(W-1)\cdots(W-r+1)=\lambda^r$.
\end{exercise}

\begin{exercise}[$\star$]\label{ex:nearbirthday}
On a circular calendar of $d=365$ days, call two birthdays \emph{close} if they are equal or adjacent. Show that the number $W$ of close pairs among $n$ people
satisfies $\|\mathcal L(W)-\text{Poisson}(3\binom n2/d)\|\le\frac{3(4n-7)}d$, and estimate the probability that $14$ people have no close pair.
\end{exercise}

\begin{exercise}[$\star\star$]\label{ex:headruns}
Flip $n+r$ fair coins $X_0,X_1,\dots,X_{n+r-1}$, and for $1\le i\le n$ let $I_i=\mathbf 1\{X_{i-1}=\mathrm T,\ X_i=\dots=X_{i+r-1}=\mathrm H\}$: a run of at least $r$ heads starts at $i$.
Show that $W=\sum_iI_i$ satisfies $\|\mathcal L(W)-\text{Poisson}(n2^{-r-1})\|\le(2r+1)2^{-r-1}$. For $n=10^6$, which $r$ makes $P(W=0)$ closest to $\frac12$, and how good is the bound there?
\end{exercise}

\begin{exercise}[$\star\star$]\label{ex:latticedw}
Let $W_n=(\varepsilon_1+\dots+\varepsilon_n)/\sqrt n$ with $\varepsilon_i=\pm1$ fair and independent. Show that $d_W(W_n,Z)\ge\frac1{2\sqrt n}-\frac{4\varphi(0)}n$, so the order of
Theorem~\ref{thm:steinclt} cannot be improved.
\end{exercise}

\begin{exercise}[$\star\star$]\label{ex:exchclt}
In the setting of Theorem~\ref{thm:steinclt}, let $X_1',\dots,X_n'$ be an independent copy of $X_1,\dots,X_n$, let $J$ be uniform on $\{1,\dots,n\}$ and independent of
everything, and $W'=W-X_J+X_J'$. Show that $(W,W')$ is exchangeable with $\mathbb E[W'-W\mid X_1,\dots,X_n]=-\frac Wn$, and deduce from Theorem~\ref{thm:exch} that
$d_W(W,Z)\le\sqrt{\sum_i\mathbb EX_i^4}+4\sum_i\mathbb E|X_i|^3$.
\end{exercise}

\section{Hints}

\begin{hint}{ex:gaussmoments}
Take $f(w)=w^{2k-1}$, after checking that Lemma~\ref{lem:stein}(i) extends to polynomials. For the second part, look at $(X-\mu)/s$.
\end{hint}
\begin{hint}{ex:dwlaw}
Approximate the bounded continuous test functions of convergence in law by Lipschitz ones; or use $h(x)=\cos tx$ and $\sin tx$, which are $|t|$-Lipschitz.
\end{hint}
\begin{hint}{ex:poissonchar}
For the converse take $g=\mathbf 1_{\{k\}}$. For the factorial moments take $g(w)=w(w-1)\cdots(w-r+2)$, after a truncation argument, or compute directly.
\end{hint}
\begin{hint}{ex:nearbirthday}
Copy Example~\ref{ex:hashcs}. Given person $i$'s birthday, the events ``$j$ close to $i$'' and ``$k$ close to $i$'' are independent.
\end{hint}
\begin{hint}{ex:headruns}
Two runs cannot \emph{start} within distance $r$ of each other: if $0<j-i\le r$ then $X_{j-1}$ would have to be both heads and tails. So $b_2=0$.
\end{hint}
\begin{hint}{ex:latticedw}
$W_n$ lives on a lattice of spacing $\delta=\frac2{\sqrt n}$. Test with $h(x)=$ distance from $x$ to that lattice.
\end{hint}
\begin{hint}{ex:exchclt}
Compute $\mathbb E[X_J'-X_J\mid X]$ and $\mathbb E[(X_J'-X_J)^2\mid X]$, and use Theorem~\ref{thm:exch} with $\mathcal G=\sigma(X_1,\dots,X_n)$. For the third moment, $|a-b|^3\le4(|a|^3+|b|^3)$.
\end{hint}

\section{Solutions}

\begin{solution}{ex:gaussmoments}
\idea The identity $\mathbb E\,Zf(Z)=\mathbb Ef'(Z)$ lowers the power of $Z$ by two at the cost of a factor: a recursion for the moments.
\writeup Lemma~\ref{lem:stein}(i) holds for any $C^1$ function $f$ of polynomial growth with $f'$ of polynomial growth: the proof only used that the boundary terms
$f(\pm R)\varphi(\pm R)$ vanish and the integrals converge. With $f(w)=w^{2k-1}$, $\mathbb EZ^{2k}=(2k-1)\mathbb EZ^{2k-2}$, so $\mathbb EZ^{2k}=(2k-1)(2k-3)\cdots1$.
For the second part, $X\sim N(\mu,s^2)$ if and only if $Y=(X-\mu)/s\sim N(0,1)$. Given $f$, put $f_1(y)=f(\mu+sy)$, so $f_1'(y)=sf'(\mu+sy)$; the condition
$\mathbb E(X-\mu)f(X)=s^2\mathbb Ef'(X)$ reads $s\,\mathbb E\,Yf_1(Y)=s\,\mathbb Ef_1'(Y)$, and as $f$ runs over the bounded $C^1$ functions with bounded derivative so does $f_1$.
Lemma~\ref{lem:stein} concludes.
\end{solution}

\begin{solution}{ex:dwlaw}
\idea Convergence in law can be tested on characteristic functions, and $\cos tx$, $\sin tx$ are Lipschitz.
\writeup For fixed $t$, $x\mapsto\cos tx$ and $x\mapsto\sin tx$ are $|t|$-Lipschitz, so $|\mathbb Ee^{itX_n}-\mathbb Ee^{itY}|\le2|t|\,d_W(X_n,Y)\to0$. L\'evy's continuity theorem
(Theorem~\ref{thm:levycont}) gives $X_n\Rightarrow Y$.
\end{solution}

\begin{solution}{ex:poissonchar}
\idea Testing the identity on indicators of single points gives a recursion between consecutive probabilities: the recursion of the Poisson law.
\writeup The direct implication is the computation before~\eqref{eq:chensteineq}. Conversely, take $g=\mathbf 1_{\{k\}}$ with $k\ge1$: $kP(W=k)=\lambda P(W=k-1)$. So
$P(W=k)=\frac{\lambda^k}{k!}P(W=0)$, and the probabilities add up to $1$ only if $P(W=0)=e^{-\lambda}$. For the factorial moments of $W\sim$ Poisson$(\lambda)$, compute
directly: $\mathbb E\,W(W-1)\cdots(W-r+1)=\sum_{k\ge r}\frac{k!}{(k-r)!}e^{-\lambda}\frac{\lambda^k}{k!}=\lambda^r\sum_{k\ge r}e^{-\lambda}\frac{\lambda^{k-r}}{(k-r)!}=\lambda^r$. (The identity with
$g(w)=w(w-1)\cdots(w-r+2)$, justified by truncating $g$ and monotone convergence, gives the same recursion $\mathbb E(W)_r=\lambda\,\mathbb E(W+1)_{r-1}$.)
\end{solution}

\begin{solution}{ex:nearbirthday}
\idea Same neighbourhoods as for exact birthdays; overlapping pairs are again independent because they are independent given the shared person.
\writeup For a pair $\{i,j\}$, $P(\text{close})=\frac3d$, since $j$'s birthday must be one of three days. The indicator depends only on the birthdays of $i$ and $j$, so take
$B_\alpha$ to be the $2n-3$ pairs meeting $\alpha$. For $\alpha=\{i,j\}$, $\beta=\{i,k\}$, given $i$'s birthday the two events are independent, each of probability $\frac3d$,
so $\mathbb E\,I_\alpha I_\beta=\frac9{d^2}$. Then $b_1+b_2=\binom n2(4n-7)\frac9{d^2}$, $\lambda=\frac{3\binom n2}d$ and the bound is $(4n-7)\frac3d(1-e^{-\lambda})\le\frac{3(4n-7)}d$.
For $n=14$, $\lambda=\frac{273}{365}=0.748$, $e^{-\lambda}=0.473$, with the bound $0.21$ on the error. The exact value, by counting arrangements on the circle, is
$0.463$: fourteen people suffice for even odds of two birthdays at most a day apart, against $23$ for an exact match (Diaconis and Mosteller, who credit the
result to Abramson and Moser).
\end{solution}

\begin{solution}{ex:headruns}
\idea Counting only the \emph{starts} of runs, marked by a tail just before, removes the clumping: a long run of heads contains many runs of $r$ heads, but only
one start.
\writeup $p_i=2^{-r-1}$ and $\lambda=n2^{-r-1}$. $I_i$ depends on $X_{i-1},\dots,X_{i+r-1}$, so it is independent of the $I_j$ with $|j-i|>r$; let $B_i=\{j:|j-i|\le r\}$,
of size at most $2r+1$. If $0<j-i\le r$ then $i\le j-1\le i+r-1$, so $I_i=1$ forces $X_{j-1}=\mathrm H$ and $I_j=1$ forces $X_{j-1}=\mathrm T$: $I_iI_j=0$ and $b_2=0$.
Then $b_1\le n(2r+1)2^{-2r-2}$, and $\frac{b_1}\lambda\le(2r+1)2^{-r-1}$. For $n=10^6$, $P(W=0)\approx e^{-\lambda}$ is closest to $\frac12$ when $\lambda=n2^{-r-1}$ is
closest to $\log2$, that is $2^{r+1}\approx1.44\times10^6$, $r=19$, with $\lambda=0.95$, $e^{-\lambda}=0.39$ and error at most $39\cdot2^{-20}<4\times10^{-5}$.
So a million flips contain a run of $19$ heads with probability about $0.61$, and that figure is guaranteed to within $4\times10^{-5}$.
(Runs starting at $X_0$ are not counted; they add at most $2^{-r}$ to the probability of a long run.)
\end{solution}

\begin{solution}{ex:latticedw}
\idea A lattice law is far from a density at the scale of the lattice: the distance-to-the-lattice function vanishes on one and averages a quarter of the
spacing on the other.
\writeup $W_n$ takes values in $L=\{\frac{n-2k}{\sqrt n}\}$, a subset of a lattice $L'=x_0+\delta\Z$ with $\delta=\frac2{\sqrt n}$. Let $h(x)=\operatorname{dist}(x,L')$, which is
$1$-Lipschitz, periodic of period $\delta$, with values in $[0,\frac\delta2]$ and mean value $\frac\delta4$ over a period. Then $\mathbb Eh(W_n)=0$. On a period cell
$C$, $\varphi$ varies by at most $\int_C|\varphi'|$, so
\[
\Bigl|\int_Ch\varphi-\frac\delta4\int_C\varphi\Bigr|\le\frac\delta2\cdot\delta\int_C|\varphi'| ,
\]
because $\int_C(h-\frac\delta4)=0$ lets us replace $\varphi$ by $\varphi-\varphi(c)$ for any $c\in C$. Summing over cells, $\mathbb Eh(Z)\ge\frac\delta4-\frac{\delta^2}2\int|\varphi'|=\frac\delta4-\delta^2\varphi(0)$,
as $\int|\varphi'|=2\varphi(0)$. So $d_W\ge\frac1{2\sqrt n}-\frac{4\varphi(0)}n$. The upper bound of Theorem~\ref{thm:steinclt} is $\frac3{\sqrt n}$.
\end{solution}

\begin{solution}{ex:exchclt}
\idea Replacing one coordinate by a fresh copy is a reversible move; on average it shrinks $W$ by $W/n$, which is the linear regression the theorem asks for.
\writeup Given $J=j$, $(X,X')$ with coordinate $j$ exchanged has the same law as $(X,X')$, so $(W,W')\overset{\mathrm d}=(W',W)$. Next
$\mathbb E[X_J'-X_J\mid X]=\frac1n\sum_j(0-X_j)=-\frac Wn$, so $\lambda=\frac1n$, and $\mathbb E[D^2\mid X]=\frac1n\sum_j(X_j^2+\sigma_j^2)$. Hence
\[
\mathbb E\Bigl|1-\frac{\mathbb E[D^2\mid X]}{2\lambda}\Bigr|=\frac12\mathbb E\Bigl|\sum_j(X_j^2-\sigma_j^2)\Bigr|\le\frac12\Bigl(\sum_j\operatorname{Var}X_j^2\Bigr)^{1/2}\le\frac12\Bigl(\sum_j\mathbb EX_j^4\Bigr)^{1/2},
\]
by Cauchy--Schwarz and independence. And $\mathbb E|D|^3=\frac1n\sum_j\mathbb E|X_j'-X_j|^3\le\frac8n\sum_j\mathbb E|X_j|^3$, so $\frac{\mathbb E|D|^3}{2\lambda}\le4\sum_j\mathbb E|X_j|^3$. Theorem~\ref{thm:exch}
gives the bound. It is weaker than Theorem~\ref{thm:steinclt}, which needs no fourth moments, but its proof never looked inside $W$: only at how $W$ moves.
\end{solution}

\section*{Sources and further reading}

\begin{itemize}
\item \emph{History.} C.~Stein, ``Inadmissibility of the usual estimator for the mean of a multivariate normal distribution'', \emph{Proc.\ Third Berkeley Symp.\ Math.\ Statist.\
Prob.}\ 1 (1956), 197--206. C.~Stein, ``A bound for the error in the normal approximation to the distribution of a sum of dependent random variables'',
\emph{Proc.\ Sixth Berkeley Symp.\ Math.\ Statist.\ Prob.}\ 2 (1972), 583--602 (the symposium was held in 1970). L.~H.~Y.~Chen, \emph{Poisson Approximation for Sums of Dependent Bernoulli
Random Variables}, PhD thesis, Stanford University, 1972, advised by Stein (Mathematics Genealogy Project), and ``Poisson approximation for dependent trials'',
\emph{Ann.\ Probab.}\ 3 (1975), 534--545. R.~Arratia, L.~Goldstein and L.~Gordon, ``Two moments suffice for Poisson approximations: the Chen--Stein method'',
\emph{Ann.\ Probab.}\ 17 (1989), 9--25. L.~H.~Y.~Chen, ``Stein's method of normal approximation: some recollections and reflections'', arXiv:2104.08302 (2021),
for Stein's discovery of the method in the words of his student.
\item C.~Stein, \emph{Approximate Computation of Expectations}, IMS Lecture Notes--Monograph Series~7, 1986. The method by its author.
\item L.~H.~Y.~Chen, L.~Goldstein and Q.-M.~Shao, \emph{Normal Approximation by Stein's Method}, Springer, 2011. Lemma~2.4 has the sharp bounds on the solution;
Chapter~3 gets Berry--Esseen with the optimal rate.
\item A.~D.~Barbour, L.~Holst and S.~Janson, \emph{Poisson Approximation}, Oxford University Press, 1992. Lemma~1.1.1 is Lemma~\ref{lem:chensteinsol}; the book treats
couplings, random graphs and much more.
\item N.~Ross, ``Fundamentals of Stein's method'', \emph{Probab.\ Surveys} 8 (2011), 210--293. The best place to start.
\item P.~Diaconis and F.~Mosteller, ``Methods for studying coincidences'', \emph{J.\ Amer.\ Statist.\ Assoc.}\ 84 (1989), 853--861. Birthdays, near-birthdays and why
coincidences are not surprising.
\end{itemize}
